\begin{table*}[!th]
	\centering
	\resizebox{1.0\linewidth}{!}{
		\begin{tabular}{ l | c c c | c c c c c c c c c c c c c}
			\toprule[1pt]
			Method & OA & mAcc & mIoU & ceiling & floor & wall & beam & column & window & door & table & chair & sofa & bookcase & board & clutter \\
			\hline
			PointNet~\cite{qi2017pointnet} & -- & 49.0 & 41.1 & 88.8 & 97.3 & 69.8 & 0.1 & 3.9 & 46.3 & 10.8 & 59.0 & 52.6 & 5.9 & 40.3 & 26.4 & 33.2\\
			SegCloud~\cite{tchapmi2017segcloud} & -- & 57.4 & 48.9 & 90.1 & 96.1 & 69.9 & 0.0 & 18.4 & 38.4 & 23.1 & 70.4 & 75.9 & 40.9 & 58.4 & 13.0 & 41.6\\
			TangentConv~\cite{tatarchenko2018tangent} & -- & 62.2 & 52.6 & 90.5 & 97.7 & 74.0 & 0.0 & 20.7 & 39.0 & 31.3 & 77.5 & 69.4 & 57.3 & 38.5 & 48.8 & 39.8 \\
			PointCNN~\cite{li2018pointcnn} & 85.9 & 63.9 & 57.3 & 92.3 & 98.2 & 79.4 & 0.0 & 17.6 & 22.8 & 62.1 & 74.4 & 80.6 & 31.7 & 66.7 & 62.1 & 56.7\\
			SPGraph~\cite{landrieu2018spg} & 86.4 & 66.5 & 58.0 & 89.4 & 96.9 & 78.1 & 0.0 & 42.8 & 48.9 & 61.6 & 84.7 & 75.4 & 69.8 & 52.6 & 2.1 & 52.2\\
			PCCN~\cite{wang2018pccn} & -- & 67.0 & 58.3 & 92.3 & 96.2 & 75.9 & 0.3 & 6.0 & 69.5 & 63.5 & 66.9 & 65.6 & 47.3 & 68.9 & 59.1 & 46.2\\
			PAT~\cite{yang2019modeling} & -- & 70.8 & 60.1 & 93.0 & 98.5 & 72.3 & 1.0 & 41.5 & 85.1 & 38.2 & 57.7 & 83.6 & 48.1 & 67.0 & 61.3 & 33.6\\
			PointWeb~\cite{zhao2019pointweb} & 87.0 & 66.6 & 60.3 & 92.0 & 98.5 & 79.4 & 0.0 & 21.1 & 59.7 & 34.8 & 76.3 & 88.3 & 46.9 & 69.3 & 64.9 & 52.5 \\
			HPEIN~\cite{jiang2019hpein} & 87.2 & 68.3 & 61.9 & 91.5 & 98.2 & 81.4 & 0.0 & 23.3 & 65.3 & 40.0 & 75.5 & 87.7 & 58.5 & 67.8 & 65.6 & 49.4 \\
			MinkowskiNet~\cite{thomas2019kpconv} & -- & 71.7 & 65.4 & 91.8 & 98.7 & 86.2 & 0.0 & 34.1 & 48.9 & 62.4 & 81.6 & 89.8 & 47.2 & 74.9 & 74.4 & 58.6 \\
			KPConv~\cite{thomas2019kpconv} & -- & 72.8 & 67.1 &92.8 & 97.3 & 82.4 & 0.0 & 23.9 & 58.0 & 69.0 & 81.5 & 91.0 & 75.4 & 75.3 & 66.7 & 58.9 \\
			\hline
			PointTransformer & \textbf{90.8} & \textbf{76.5} & \textbf{70.4} & 94.0 & 98.5 & 86.3 & 0.0 & 38.0 & 63.4 & 74.3 & 89.1 & 82.4 & 74.3 & 80.2 & 76.0 & 59.3 \\
			\bottomrule[1pt]
	\end{tabular}}
	\caption{Semantic segmentation results on the S3DIS dataset, evaluated on Area 5.}
	\label{tab:s3disresult}
\end{table*}

\begin{table}[!th]
	\centering
		\begin{tabular}{ l | c c c}
			\toprule[1pt]
			Method & OA & mAcc & mIoU \\
			\hline
			PointNet~\cite{qi2017pointnet} & 78.5 & 66.2 & 47.6 \\
			RSNet~\cite{huang2018recurrent} & -- & 66.5 & 56.5 \\
			SPGraph~\cite{landrieu2018spg} & 85.5 & 73.0 & 62.1 \\
			PAT~\cite{yang2019modeling} & -- & 76.5 & 64.3 \\
			PointCNN~\cite{li2018pointcnn} & 88.1 & 75.6 & 65.4 \\
			PointWeb~\cite{zhao2019pointweb} & 87.3 & 76.2 & 66.7 \\
			ShellNet~\cite{zhang2019shellnet} & 87.1 & -- & 66.8 \\
			RandLA-Net~\cite{thomas2019kpconv} & 88.0 & 82.0 & 70.0 \\
			KPConv~\cite{thomas2019kpconv} & -- & 79.1 & 70.6 \\
			\hline
			PointTransformer & \textbf{90.2} & \textbf{81.9} & \textbf{73.5} \\
			\bottomrule[1pt]
	\end{tabular}
	\caption{Semantic segmentation results on the S3DIS dataset, evaluated with 6-fold cross-validation.}
	\label{tab:s3disresult2}
\end{table}

We evaluate the effectiveness of the presented Point Transformer design on a number of domains and tasks. For 3D semantic segmentation, we use the challenging Stanford Large-Scale 3D Indoor Spaces (S3DIS) dataset~\cite{armeni2016s3dis}. For 3D shape classification, we use the widely adopted ModelNet40 dataset~\cite{wu2015modelnet}. And for object part segmentation, we use ShapeNetPart~\cite{yi2016scalable}.

\mypara{Implementation details.}
We implement the Point Transformer in PyTorch~\cite{Paszke2019Pytorch}. We use the SGD optimizer with momentum and weight decay set to 0.9 and 0.0001, respectively. For semantic segmentation on S3DIS, we train for 40K iterations with initial learning rate 0.5, dropped by 10x at steps 24K and 32K. For 3D shape classification on ModelNet40 and 3D object part segmentation on ShapeNetPart, we train for 200 epochs. The initial learning rate is set to 0.05 and is dropped by 10x at epochs 120 and 160.

\subsection{Semantic Segmentation}

\noindent\textbf{Data and metric.}
The S3DIS~\cite{armeni2016s3dis} dataset for semantic scene parsing consists of 271 rooms in six areas from three different buildings. Each point in the scan is assigned a semantic label from 13 categories (ceiling, floor, table, etc.). Following a common protocol~\cite{tchapmi2017segcloud,qi2017pointnet2}, we evaluate the presented approach in two modes: (a) Area 5 is withheld during training and is used for testing, and (b) 6-fold cross-validation. For evaluation metrics, we use mean classwise intersection over union (mIoU), mean of classwise accuracy (mAcc), and overall pointwise accuracy (OA).

\mypara{Performance comparison.}
The results are presented in Tables~\ref{tab:s3disresult} and~\ref{tab:s3disresult2}. The Point Transformer outperforms all prior models according to all metrics in both evaluation modes. On Area 5, the Point Transformer attains mIoU/mAcc/OA of 70.4\%/76.5\%/90.8\%, outperforming all prior work by multiple percentage points in each metric. The Point Transformer is the first model to pass the 70\% mIoU bar, outperforming the prior state of the art by 3.3 absolute percentage points in mIoU. The Point Transformer outperforms MLPs-based frameworks such as PointNet~\cite{qi2017pointnet}, voxel-based architectures such as SegCloud~\cite{tchapmi2017segcloud}, graph-based methods such as SPGraph~\cite{landrieu2018spg}, attention-based methods such as PAT~\cite{yang2019modeling}, sparse convolutional networks such as MinkowskiNet~\cite{choy20194d}, and continuous convolutional networks such as KPConv~\cite{thomas2019kpconv}.
Point Transformer also substantially outperforms all prior models under 6-fold cross-validation. The mIoU in this mode is 73.5\%, outperforming the prior state of the art (KPConv) by 2.9 absolute percentage points. The number of parameters in Point Transformer (4.9M) is much smaller than in current high-performing architectures such as KPConv (14.9M) and SparseConv (30.1M).

\mypara{Visualization.}
Figure~\ref{fig:s3disvisual} shows the Point Transformer's predictions.
We can see that the predictions are very close to the ground truth. Point Transformer captures detailed semantic structure in complex 3D scenes, such as the legs of chairs, the outlines of poster boards, and the trim around doorways.

\begin{table}[h]
	\centering
	\begin{tabular}{ l | c c c}
		\toprule[1pt]
		Method & input & mAcc & OA \\
		\hline
		3DShapeNets~\cite{wu2015modelnet} & voxel & 77.3 & 84.7 \\
		VoxNet~\cite{maturana2015voxnet} & voxel & 83.0 & 85.9 \\
		Subvolume~\cite{qi2016volumetric} & voxel & 86.0 & 89.2 \\
		MVCNN~\cite{su15mvcnn} & image & -- & 90.1 \\
		PointNet~\cite{qi2017pointnet} & point & 86.2 & 89.2 \\
		A-SCN~\cite{xie2018attentional} & point & 87.6 & 90.0 \\
		Set Transformer~\cite{lee2019set} & point & -- & 90.4 \\
		PAT~\cite{yang2019modeling} & point & -- & 91.7 \\
		PointNet++~\cite{qi2017pointnet2} & point & -- & 91.9 \\
		SpecGCN~\cite{wang2018local} & point & -- & 92.1 \\
		PointCNN~\cite{li2018pointcnn} & point & 88.1 & 92.2 \\
		DGCNN~\cite{wang2019dgcnn} & point & 90.2 & 92.2 \\
		PointWeb~\cite{zhao2019pointweb} & point & 89.4 & 92.3 \\
		SpiderCNN~\cite{xu2018spidercnn} & point & -- & 92.4 \\
		PointConv~\cite{wu2019pointconv} & point & -- & 92.5 \\
		Point2Sequence~\cite{liu2019point2sequence} & point & 90.4 & 92.6 \\
		KPConv~\cite{thomas2019kpconv} & point & -- & 92.9 \\
		InterpCNN~\cite{mao2019interpolated} & point & -- & 93.0 \\
		\hline
		PointTransformer & point & \textbf{90.6} & \textbf{93.7} \\
		\bottomrule[1pt]
	\end{tabular}
	\caption{Shape classification results on the ModelNet40 dataset.}
	\label{tab:modelnet40result}
	\vspace{-2mm}
\end{table}

\subsection{Shape Classification}

\noindent\textbf{Data and metric.}
The ModelNet40~\cite{wu2015modelnet} dataset contains 12,311 CAD models with 40 object categories. They are split into 9,843 models for training and 2,468 for testing. We follow the data preparation procedure of Qi et al.~\cite{qi2017pointnet2} and uniformly sample the points from each CAD model together with the normal vectors from the object meshes. For evaluation metrics, we use the mean accuracy within each category (mAcc) and the overall accuracy (OA) over all classes.

\mypara{Performance comparison.}
The results are presented in Table~\ref{tab:modelnet40result}. The Point Transformer sets the new state of the art in both metrics. The overall accuracy of Point Transformer on ModelNet40 is 93.7\%. It outperforms strong graph-based models such as DGCNN~\cite{wang2019dgcnn}, attention-based models such as A-SCN~\cite{xie2018attentional} and Point2Sequence~\cite{liu2019point2sequence},
and strong point-based models such as KPConv~\cite{thomas2019kpconv}.

\mypara{Visualization.}
To probe the representation learned by the Point Transformer, we conduct shape retrieval by retrieving nearest neighbors in the space of the output features produced by the Point Transformer. Some results are shown in Figure~\ref{fig:modelnet40visual}. The retrieved shapes are very similar to the query, and when they differ, they differ along aspects that we perceive as less semantically salient, such as legs of desks.

\begin{table}[t]
	\centering
	\begin{tabular}{ l | c c}
		\toprule[1pt]
		Method & cat.\ mIoU & ins.\ mIoU \\
		\hline
		PointNet~\cite{qi2017pointnet} & 80.4 & 83.7 \\
		A-SCN~\cite{xie2018attentional} & -- & 84.6 \\
		PCNN~\cite{wang2018pccn} & 81.8 & 85.1 \\
		PointNet++~\cite{qi2017pointnet2} & 81.9 & 85.1 \\
		DGCNN~\cite{wang2019dgcnn} & 82.3 & 85.1 \\
		Point2Sequence~\cite{liu2019point2sequence} & -- & 85.2 \\
		SpiderCNN~\cite{xu2018spidercnn}& 81.7 & 85.3 \\
		SPLATNet~\cite{su2018splatnet} & 83.7 & 85.4 \\
		PointConv~\cite{wu2019pointconv} & 82.8 & 85.7 \\
		SGPN~\cite{wang2018sgpn} & 82.8 & 85.8 \\
		PointCNN~\cite{li2018pointcnn} & 84.6 & 86.1 \\
		InterpCNN~\cite{mao2019interpolated} & 84.0 & 86.3 \\
		KPConv~\cite{thomas2019kpconv} & \textbf{85.1} & 86.4 \\
		\hline
		PointTransformer & 83.7 & \textbf{86.6} \\
		\bottomrule[1pt]
	\end{tabular}
	\caption{Object part segmentation results on the ShapeNetPart dataset.}
	\label{tab:shapenetpartresult}
	\vspace{-2mm}
\end{table}

\subsection{Object Part Segmentation}

\noindent\textbf{Data and metric.}
The ShapeNetPart dataset~\cite{yi2016scalable} is annotated for 3D object part segmentation. It consists of 16,880 models from 16 shape categories, with 14,006 3D models for training and 2,874 for testing. The number of parts for each category is between 2 and 6, with 50 different parts in total.
We use the sampled point sets produced by Qi et al.~\cite{qi2017pointnet2} for a fair comparison with prior work. For evaluation metrics, we report category mIoU and instance mIoU.

\mypara{Performance comparison.}
The results are presented in Table~\ref{tab:shapenetpartresult}. The Point Transformer outperforms all prior models as measured by instance mIoU.
(Note that we did not use loss-balancing during training, which can boost category mIoU.)

\mypara{Visualization.}
Object part segmentation results on a number of models are shown in Figure~\ref{fig:shapenetpartvisual}. The Point Transformer's part segmentation predictions are clean and close to the ground truth.

\begin{figure*}
	\centering
	\small
	\resizebox{0.97\linewidth}{!}{
	\begin{tabular}{@{\hspace{0.0mm}}c@{\hspace{1.0mm}}c@{\hspace{1.0mm}}c@{\hspace{2.5mm}}c@{\hspace{1.0mm}}c@{\hspace{1.0mm}}c@{\hspace{0.0mm}}}
		Input & Ground Truth & Point Transformer & Input & Ground Truth & Point Transformer \\
		\includegraphics[width=0.15\linewidth]{figure/s3dis/conference_00.jpg}&
		\includegraphics[width=0.15\linewidth]{figure/s3dis/conference_01.jpg}&
		\includegraphics[width=0.15\linewidth]{figure/s3dis/conference_02.jpg}&
		\includegraphics[width=0.15\linewidth]{figure/s3dis/lobby1_00.jpg}&
		\includegraphics[width=0.15\linewidth]{figure/s3dis/lobby1_01.jpg}&
		\includegraphics[width=0.15\linewidth]{figure/s3dis/lobby1_02.jpg}\\
		\includegraphics[width=0.15\linewidth]{figure/s3dis/office6_00.jpg}&
		\includegraphics[width=0.15\linewidth]{figure/s3dis/office6_01.jpg}&
		\includegraphics[width=0.15\linewidth]{figure/s3dis/office6_02.jpg}&
		\includegraphics[width=0.15\linewidth]{figure/s3dis/hallway5_00.jpg}&
		\includegraphics[width=0.15\linewidth]{figure/s3dis/hallway5_01.jpg}&
		\includegraphics[width=0.15\linewidth]{figure/s3dis/hallway5_02.jpg}\\
		\includegraphics[width=0.15\linewidth]{figure/s3dis/hallway_00.jpg}&
		\includegraphics[width=0.15\linewidth]{figure/s3dis/hallway_01.jpg}&
		\includegraphics[width=0.15\linewidth]{figure/s3dis/hallway_02.jpg}&
		\includegraphics[width=0.15\linewidth]{figure/s3dis/office42_00.jpg}&
		\includegraphics[width=0.15\linewidth]{figure/s3dis/office42_01.jpg}&
		\includegraphics[width=0.15\linewidth]{figure/s3dis/office42_02.jpg}\\
		\multicolumn{6}{c}{\includegraphics[width=0.9\linewidth]{figure/colorbar.pdf}}\\
	\end{tabular}}
	\caption{Visualization of semantic segmentation results on the S3DIS dataset.}
	\label{fig:s3disvisual}
\end{figure*}

\begin{figure*}
	\centering
	\resizebox{0.97\linewidth}{!}{
	\begin{tabular}{c|cccccc}
	\includegraphics[width=0.16\linewidth]{figure/modelnet40/car/snapshot00.png}&
	\includegraphics[width=0.16\linewidth]{figure/modelnet40/car/snapshot01.png}&
	\includegraphics[width=0.16\linewidth]{figure/modelnet40/car/snapshot02.png}&
	\includegraphics[width=0.16\linewidth]{figure/modelnet40/car/snapshot03.png}&
	\includegraphics[width=0.16\linewidth]{figure/modelnet40/car/snapshot04.png}&
	\includegraphics[width=0.16\linewidth]{figure/modelnet40/car/snapshot05.png}&
	\includegraphics[width=0.16\linewidth]{figure/modelnet40/car/snapshot06.png}\\
	\includegraphics[width=0.16\linewidth]{figure/modelnet40/desk/snapshot00.png}&
	\includegraphics[width=0.16\linewidth]{figure/modelnet40/desk/snapshot01.png}&
	\includegraphics[width=0.16\linewidth]{figure/modelnet40/desk/snapshot02.png}&
	\includegraphics[width=0.16\linewidth]{figure/modelnet40/desk/snapshot03.png}&
	\includegraphics[width=0.16\linewidth]{figure/modelnet40/desk/snapshot04.png}&
	\includegraphics[width=0.16\linewidth]{figure/modelnet40/desk/snapshot05.png}&
	\includegraphics[width=0.16\linewidth]{figure/modelnet40/desk/snapshot06.png}\\
	\end{tabular}}
	\caption{Visualization of shape retrieval results on the ModelNet40 dataset. The leftmost column shows the input query and the other columns show the retrieved models.}
	\label{fig:modelnet40visual}
\end{figure*}

\begin{figure*}[!h]
	\centering
	\resizebox{0.97\linewidth}{!}{
	\begin{tabular}{ccccccccc}
	\includegraphics[width=0.195\linewidth]{figure/shapenetpart/chair_00.png}&
	\includegraphics[width=0.195\linewidth]{figure/shapenetpart/airplane_00.png}&
	\includegraphics[width=0.195\linewidth]{figure/shapenetpart/lamp_00.png}&
	\includegraphics[width=0.195\linewidth]{figure/shapenetpart/skateboard_00.png}&
	\includegraphics[width=0.195\linewidth]{figure/shapenetpart/table_00.png}&
	\includegraphics[width=0.195\linewidth]{figure/shapenetpart/gun_00.png}&
	\includegraphics[width=0.195\linewidth]{figure/shapenetpart/knife_00.png}&
	\includegraphics[width=0.195\linewidth]{figure/shapenetpart/car_00.png}&
	\includegraphics[width=0.195\linewidth]{figure/shapenetpart/guitar_00.png}\\
	\includegraphics[width=0.195\linewidth]{figure/shapenetpart/chair_01.png}&
	\includegraphics[width=0.195\linewidth]{figure/shapenetpart/airplane_01.png}&
	\includegraphics[width=0.195\linewidth]{figure/shapenetpart/lamp_01.png}&
	\includegraphics[width=0.195\linewidth]{figure/shapenetpart/skateboard_01.png}&
	\includegraphics[width=0.195\linewidth]{figure/shapenetpart/table_01.png}&
	\includegraphics[width=0.195\linewidth]{figure/shapenetpart/gun_01.png}&
	\includegraphics[width=0.195\linewidth]{figure/shapenetpart/knife_01.png}&
	\includegraphics[width=0.195\linewidth]{figure/shapenetpart/car_01.png}&
	\includegraphics[width=0.195\linewidth]{figure/shapenetpart/guitar_01.png}\\
	\end{tabular}}
	\caption{Visualization of object part segmentation results on the ShapeNetPart dataset. The ground truth is in the top row, Point Transformer predictions on the bottom.}
	\label{fig:shapenetpartvisual}
\end{figure*}

\begin{table}
	\centering
		\begin{tabular}{ c | c c c }
			\toprule[1pt]
			$k$ & mIoU & mAcc & OA \\
			\hline
			4 & 59.6 & 66.0 & 86.0 \\
			8 & 67.7 & 73.8 & 89.9 \\
			16 & \textbf{70.4} & \textbf{76.5} & \textbf{90.8} \\
			32 & 68.3 & 75.0 & 89.8 \\
			64 & 67.7 & 74.1 & 89.9 \\
			\bottomrule[1pt]
		\end{tabular}
	\caption{Ablation study: number of neighbors $k$ in the definition of local neighborhoods.}
	\label{tab:s3disknn}
\end{table}

\begin{table}
	\centering
		\begin{tabular}{ c | c c c }
			\toprule[1pt]
			Pos.\ encoding & mIoU & mAcc & OA \\
			\hline
			none & 64.6 & 71.9 & 88.2 \\
			absolute & 66.5 & 73.2 & 88.9 \\
			relative & \textbf{70.4} & \textbf{76.5} & \textbf{90.8} \\
			relative for attention & 67.0 & 73.0 & 89.3 \\
			relative for feature & 68.7 & 74.4 & 90.4 \\
			\bottomrule[1pt]
		\end{tabular}
	\caption{Ablation study: position encoding.}
	\label{tab:s3disposition}
\end{table}

\begin{table}
	\centering
		\begin{tabular}{ c | c c c }
			\toprule[1pt]
			Operator & mIoU & mAcc & OA \\
			\hline
			MLP & 61.7 & 68.6 & 87.1 \\
			MLP+pooling & 63.7 & 71.0 & 87.8 \\
			\hline
			scalar attention & 64.6 & 71.9 & 88.4 \\
			vector attention & \textbf{70.4} & \textbf{76.5} & \textbf{90.8} \\
			\bottomrule[1pt]
		\end{tabular}
	\caption{Ablation study: form of self-attention operator.}
	\label{tab:s3disattention}
\end{table}

\subsection{Ablation Study}
\label{sec:ablation}

We now conduct a number of controlled experiments that examine specific decisions in the Point Transformer design. These studies are performed on the semantic segmentation task on the S3DIS dataset, tested on Area 5.

\mypara{Number of neighbors.}
We first investigate the setting of the number of neighbors $k$, which is used in determining the local neighborhood around each point. The results are shown in Table~\ref{tab:s3disknn}. The best performance is achieved when $k$ is set to 16. When the neighborhood is smaller ($k=4$ or $k=8$), the model may not have sufficient context for its predictions. When the neighborhood is larger ($k=32$ or $k=64$), each self-attention layer is provided with a large number of datapoints, many of which may be farther and less relevant. This may introduce excessive noise into the processing, lowering the model's accuracy.

\mypara{Softmax regularization.}
We conduct an ablation study on the normalization function $\rho$ in Eq.~\ref{eq:pointtransformer}. The performance without softmax regularization on S3DIS Area5 is 66.5\%/72.8\%/89.3\%, in terms of mIoU/mAcc/OA. It is much lower than the performance with softmax regularization (70.4\%/76.5\%90.8\%). This suggests that the normalization is essential in this setting.

\mypara{Position encoding.}
We now study the choice of the position encoding $\delta$. The results are shown in Table~\ref{tab:s3disposition}. We can see that without position encoding, the performance drops significantly. With absolute position encoding, the performance is higher than without. Relative position encoding yields the highest performance. When relative position encoding is added only to the attention generation branch (first term in Eq.~\ref{eq:pointtransformer}) or only to the feature transformation branch (second term in Eq.~\ref{eq:pointtransformer}), the performance drops again, indicating that adding the relative position encoding to both branches is important.

\mypara{Attention type.}
Finally, we investigate the type of self-attention used in the point transformer layer. The results are shown in Table~\ref{tab:s3disattention}. We examine four conditions. `MLP' is a no-attention baseline that replaces the point transformer layer in the point transformer block with a pointwise MLP. `MLP+pooling' is a more advanced no-attention baseline that replaces the point transformer layer with a pointwise MLP followed by max pooling within each $k$NN neighborhood: this performs feature transformation at each point and enables each point to exchange information with its local neighborhood, but does not leverage attention mechanisms.
`scalar attention' replaces the vector attention used in Eq.~\ref{eq:pointtransformer} by scalar attention, as in Eq.~\ref{eq:scalar} and in the original transformer design~\cite{vaswani2017attention}. `vector attention' is the formulation we use, presented in Eq.~\ref{eq:pointtransformer}. We can see that scalar attention is more expressive than the no-attention baselines, but is in turn outperformed by vector attention. The performance gap between vector and scalar attention is significant: 70.4\% vs.\ 64.6\%, an improvement of 5.8 absolute percentage points. Vector attention is more expressive since it supports adaptive modulation of individual feature channels, not just whole feature vectors. This expressivity appears to be very beneficial in 3D data processing.